  \begin{block}{Key Findings}

    \begin{itemize}
      \item \textbf{50\% of} government websites embed at least one tracker by 2025, up from near zero in 1996, with a steep post-2016 surge.
      \item \textbf{Third-party services drive nearly all growth:} 48\% of sites carry third-party trackers in 2025 versus only \textbf{4.4\%} with first-party trackers.
      \item \textbf{Adoption is highly uneven:} Singapore (\textbf{95.5\%}), Australia (\textbf{89.0\%}), and Bangladesh (\textbf{82.5\%}) lead in 2025, while China (\textbf{6.1\%}), Germany (\textbf{9.2\%}), and France (\textbf{11.6\%}) stay persistently low.
      \item \textbf{Ownership is concentrated:} Google trackers were embedded \textbf{68K times}, far ahead of Meta (\textbf{6,635}); \emph{google-analytics.com} and \emph{googletagmanager.com} appear in \textbf{all 61 countries}; mean $HHI_{\mathrm{org}}$ = \textbf{0.601} vs.\ $HHI_{\mathrm{tracker}}$ = \textbf{0.374}.
      \item \textbf{Trackers sometimes vanish:} we find \textbf{39 tracker gaps} across \textbf{26 countries}; in \textbf{37} cases the same sites remained archived without trackers, suggesting deliberate temporary suspensions.
      \item \textbf{High churn:} only \textbf{30\%} of trackers persist year-over-year after first uptake, with an average of \textbf{77 new tracker uptakes per year} globally (peak \textbf{161} in 2017).
    \end{itemize}

  \end{block}
